% Lapisan solusi O001 -- Bab 5 (Operator pada Ruang Hilbert)
% 4 solusi asli terpisah; bukan tulisan John M. Erdman.
% Provenans model: OpenAI Codex gpt-5.6-sol, Ultra, atas arahan pengguna.
% Lisensi komponen solusi: CC BY-SA 4.0.
% Tidak menyiratkan dukungan John M. Erdman atau Portland State University.

\markboth{SOLUSI O001 UNTUK BAB 5: OPERATOR PADA RUANG HILBERT}{SOLUSI O001 UNTUK BAB 5: OPERATOR PADA RUANG HILBERT}
\section*{Solusi O001 untuk Bab 5: Operator pada Ruang Hilbert}
\addcontentsline{toc}{section}{Solusi O001 untuk Bab 5: Operator pada Ruang Hilbert}
\begin{center}
\small
Komponen solusi asli terpisah, CC BY-SA 4.0. Ditulis dengan bantuan
OpenAI Codex gpt-5.6-sol, Ultra, atas arahan pengguna. Pernyataan latihan
berasal dari edisi terjemahan karya John M. Erdman; solusi ini bukan tulisan
beliau dan tidak menyiratkan dukungan beliau atau Portland State University.
\end{center}

% O001-SOLUTION-ID: O001-FAOA-2015-CH05-EX-001-SOLUTION
% SOURCE-EXERCISE-ID: FAOA-2015-CH05-NODE-0034
% STATEMENT-TARGET-SHA256: ee6308dea448da22e9841970db96f86a106225688b6a5aef98ca4b39f1d20a35
\begin{o001solution}
  {O001-FAOA-2015-CH05-EX-001-SOLUTION}
  {FAOA-2015-CH05-NODE-0034}
  {ee6308dea448da22e9841970db96f86a106225688b6a5aef98ca4b39f1d20a35}
\begin{o001statement}
Misalkan $\N_3 = \{1,2,3\}$ dan $\mu$ ukuran pencacahan pada $\N_3$.
 \begin{enumerate}
  \item[(a)] Identifikasilah ruang Hilbert~$L_2(\N_3,\mu)$.
  \item[(b)] Identifikasilah operator perkalian pada~$L_2(\N_3,\mu)$.
  \item[(c)] Misalkan $T$ operator linear pada $\C^3$ yang representasi
matriksnya adalah
   \[ \begin{bmatrix}
          2  &  0  &  1 \\
          0  &  2  & -1 \\
          1  & -1  &  1
      \end{bmatrix}.\]
Gunakan operator $T$ untuk mengilustrasikan rumusan teorema spektral sebelumnya.
 \end{enumerate}
\end{o001statement}
\begin{o001answer}
(a) $L_2(\N_3,\mu)=\C^3$ dengan
$\langle x,y\rangle=\sum_{j=1}^3x_j\conj{y_j}$. (b) Operator perkalian
adalah $M_\phi=\diag(\phi(1),\phi(2),\phi(3))$. (c) Matriks $T$ mempunyai
nilai eigen $0,2,3$ dengan basis ortonormal
\[
 u_0=\frac{(-1,1,2)}{\sqrt6},\qquad
 u_2=\frac{(1,1,0)}{\sqrt2},\qquad
 u_3=\frac{(1,-1,1)}{\sqrt3}.
\]
Jadi $T$ ekuivalen secara uniter dengan operator perkalian oleh fungsi
$\phi(1)=0$, $\phi(2)=2$, $\phi(3)=3$ pada~$L_2(\N_3,\mu)$.
\end{o001answer}
\begin{o001proof}
Karena ukuran setiap singleton sama dengan satu, kelas fungsi pada $\N_3$
ditentukan oleh tiga nilainya dan
\[
 \norm f_2^2=\sum_{j=1}^3\abs{f(j)}^2.
\]
Ini memberikan identifikasi isometrik dalam (a). Perkalian titik demi titik
oleh $\phi$ kemudian mempunyai matriks diagonal yang dinyatakan dalam (b),
dan setiap operator perkalian diperoleh dengan cara itu.

Untuk matriks $A$ pada (c), ekspansi determinan memberi
\[
 \det(\lambda I-A)=\lambda(\lambda-2)(\lambda-3).
\]
Substitusi langsung menunjukkan
\[
 Au_0=0,\qquad Au_2=2u_2,\qquad Au_3=3u_3.
\]
Ketiga vektor yang ditampilkan bernorma satu dan hasil kali dalam setiap
pasangannya nol, sehingga membentuk basis ortonormal. Definisikan
$U\colon L_2(\N_3,\mu)\to\C^3$ dengan
$U\delta_1=u_0$, $U\delta_2=u_2$, dan $U\delta_3=u_3$. Maka $U$ uniter dan,
untuk $j=1,2,3$,
\[
 TU\delta_j=\phi(j)U\delta_j=UM_\phi\delta_j.
\]
Karena vektor-vektor itu suatu basis, $T=UM_\phi U^{-1}$. Ini tepat
rumusan teorema spektral: operator normal (di sini bahkan swaadjoin)
ekuivalen secara uniter dengan operator perkalian.
\end{o001proof}
\end{o001solution}

% O001-SOLUTION-ID: O001-FAOA-2015-CH05-EX-002-SOLUTION
% SOURCE-EXERCISE-ID: FAOA-2015-CH05-NODE-0038
% STATEMENT-TARGET-SHA256: a4478e29765c8fb04487e45bd6b6fd8fdd22d706e22b5fab1b6b7f7294f1952c
\begin{o001solution}
  {O001-FAOA-2015-CH05-EX-002-SOLUTION}
  {FAOA-2015-CH05-NODE-0038}
  {a4478e29765c8fb04487e45bd6b6fd8fdd22d706e22b5fab1b6b7f7294f1952c}
\begin{o001statement}
Misalkan $V$ adalah \emph{operator Volterra} yang didefinisikan di atas.
 \begin{enumerate}
  \item[(a)] Tunjukkan bahwa $V$ injektif.
  \item[(b)] Hitunglah $V^*$.
  \item[(c)] Tentukan $V + V^*$ beserta jangkauannya.
 \end{enumerate}
\end{o001statement}
\begin{o001answer}
$V$ injektif,
\[
 (V^*g)(t)=\int_t^1g(x)\,dx,
 \qquad
 ((V+V^*)f)(x)=\int_0^1f(t)\,dt,
\]
dan $\ran(V+V^*)$ adalah subruang satu-dimensi semua fungsi konstan pada
$[0,1]$.
\end{o001answer}
\begin{o001proof}
Untuk $f\in L_2([0,1])$, fungsi
$F(x)=\int_0^x f(t)\,dt$ mempunyai wakil kontinu mutlak dan $F'=f$ hampir
di mana-mana. Jika $Vf=0$ sebagai anggota $L_2$, maka $F=0$ hampir di
mana-mana. Kekontinuan $F$ memaksa $F=0$ di setiap titik, sehingga
$f=F'=0$ hampir di mana-mana. Jadi $V$ injektif.

Teorema Fubini dapat diterapkan karena, dengan Cauchy--Schwarz, integral
mutlak pada segitiga $0\leq t\leq x\leq1$ berhingga. Untuk $f,g\in L_2$,
\begin{align*}
 \langle Vf,g\rangle
 &=\int_0^1\left(\int_0^x f(t)\,dt\right)g(x)\,dx\\
 &=\int_0^1 f(t)\left(\int_t^1g(x)\,dx\right)dt.
\end{align*}
Karena ruangnya real, ini sama dengan
$\langle f,V^*g\rangle$ untuk rumus $V^*$ dalam jawaban. Akhirnya,
\[
 ((V+V^*)f)(x)
 =\int_0^xf(t)\,dt+\int_x^1f(t)\,dt
 =\int_0^1f(t)\,dt.
\]
Jadi setiap citra konstan. Sebaliknya, fungsi konstan $f=c$ menghasilkan
fungsi konstan bernilai $c$, sehingga jangkauannya tepat semua fungsi
konstan.
\end{o001proof}
\end{o001solution}

% O001-SOLUTION-ID: O001-FAOA-2015-CH05-EX-003-SOLUTION
% SOURCE-EXERCISE-ID: FAOA-2015-CH05-NODE-0051
% STATEMENT-TARGET-SHA256: a360c061f372f76160bf31105405bc47db3322f73be4c99931a6f4e167c6eb0d
\begin{o001solution}
  {O001-FAOA-2015-CH05-EX-003-SOLUTION}
  {FAOA-2015-CH05-NODE-0051}
  {a360c061f372f76160bf31105405bc47db3322f73be4c99931a6f4e167c6eb0d}
\begin{o001statement}
Misalkan $\fml E$ adalah basis ortonormal bagi ruang Hilbert~$H$. Bahas transformasi linear $T$ pada $H$
sedemikian sehingga $T(e) = 0$ untuk setiap $e \in \fml E$. Berikan contoh tak trivial.
\end{o001statement}
\begin{o001answer}
Jika $T$ terbatas, maka $T=0$. Namun pada ruang Hilbert berdimensi tak
hingga terdapat transformasi linear tak kontinu dan tak nol yang menghilang
pada setiap vektor suatu basis ortonormal.
\end{o001answer}
\begin{o001proof}
Rentang linear $\spn\fml E$ padat dalam~$H$. Jika $T$ terbatas, untuk setiap
$x\in H$ pilih jaring (atau barisan jika basisnya terhitung) $x_\alpha$ dari
$\spn\fml E$ yang konvergen dalam norma ke~$x$. Hipotesis memberi
$Tx_\alpha=0$, sehingga kekontinuan menghasilkan $Tx=0$.

Untuk contoh tak trivial, ambil $H=l_2$, $\fml E=\{e^n:n\in\N\}$, dan
\[
 u=(1,1/2,1/3,\ldots)\in l_2.
\]
Vektor $u$ tidak berada dalam rentang linear hingga $\spn\fml E$, sehingga
$\fml E\cup\{u\}$ bebas linear. Perluas himpunan ini menjadi basis Hamel
$\fml B$ untuk $l_2$. Definisikan $T$ pada basis Hamel dengan
\[
 T(u)=e^1,\qquad T(b)=0\quad(b\in\fml B\setminus\{u\}),
\]
lalu perluas secara linear. Maka $T(e^n)=0$ untuk setiap~$n$, tetapi
$T(u)=e^1\neq0$. Transformasi ini harus tak terbatas, sebab seandainya
terbatas, argumen kepadatan pada paragraf pertama akan memaksanya sama
dengan nol. Dalam dimensi hingga, basis ortonormal sekaligus basis Hamel,
sehingga contoh tak trivial tidak mungkin ada.
\end{o001proof}
\end{o001solution}

% O001-SOLUTION-ID: O001-FAOA-2015-CH05-EX-004-SOLUTION
% SOURCE-EXERCISE-ID: FAOA-2015-CH05-NODE-0075
% STATEMENT-TARGET-SHA256: 84852f969329993492bb16e9519447052b491750458f20a6faa9aa1a87069596
\begin{o001solution}
  {O001-FAOA-2015-CH05-EX-004-SOLUTION}
  {FAOA-2015-CH05-NODE-0075}
  {84852f969329993492bb16e9519447052b491750458f20a6faa9aa1a87069596}
\begin{o001statement}
Misalkan $H$ ruang Hilbert kompleks dan $T$ operator positif pada~$H$.  Definisikan
   \[ \langle x,y \rangle_1 := \langle Tx,y \rangle \qquad \text{ dan } \qquad \norm{x}_1 \equiv \sqrt{\langle x,x \rangle_1} \]
untuk semua $x \in H$.
 \begin{enumerate}
  \item[(a)] Buktikan bahwa $\norm{\hphantom O}_1$ merupakan norma jika dan hanya jika $\ker T = \{0\}$.
  \item[(b)] Misalkan $\ker T = \{0\}$. Tunjukkan bahwa norma $\norm{\hphantom O}$ dan
$\norm{\hphantom O}_1$ menginduksi topologi yang sama pada~$H$ jika dan hanya jika $T$ invertibel dalam
$\ofml B(H)$. \emph{Petunjuk.}  Ketika terdapat dua topologi pada~$H$, sering kali berguna untuk
mempertimbangkan operator identitas pada~$H$.
  \item[(c)] Misalkan $T$ invertibel dalam $\ofml B(H)$.  Untuk $S \in \ofml B(H)$, carilah suatu rumus bagi
adjoin dari $S$ terhadap hasil kali dalam $\langle \hphantom x , \hphantom y \rangle_1$.
 \end{enumerate}
\end{o001statement}
\begin{o001answer}
(a) $\norm{\hphantom O}_1$ adalah norma tepat ketika $\ker T=\{0\}$.
(b) Kedua norma menginduksi topologi yang sama tepat ketika $T$ invertibel.
(c) Jika $S^{*_1}$ menyatakan adjoin terhadap hasil kali dalam baru, maka
\[
 S^{*_1}=T^{-1}S^*T.
\]
\end{o001answer}
\begin{o001proof}
Bentuk $[x,y]=\langle Tx,y\rangle$ bersifat seskuilinear dan Hermitian
karena $T=T^*$, serta $[x,x]\geq0$. Ketaksamaan Cauchy--Schwarz untuk bentuk
positif memberi
\[
 \abs{[x,y]}^2\leq[x,x][y,y].
\]
Jika $[x,x]=0$, ruas kanan menunjukkan $\langle Tx,y\rangle=0$ untuk semua
$y$, jadi $Tx=0$. Sebaliknya, $Tx=0$ langsung memberi $[x,x]=0$. Maka
bentuk itu definit positif, dan akar kuadratnya norma, tepat ketika
$\ker T=\{0\}$. Ini membuktikan (a), termasuk ketaksamaan segitiga yang
merupakan konsekuensi Cauchy--Schwarz bagi hasil kali dalam~$[\ ,\ ]$.

Selalu berlaku
\[
 \norm{x}_1^2=\langle Tx,x\rangle\leq\norm T\,\norm x^2,
\]
jadi identitas $I\colon(H,\norm{\hphantom O})\to
(H,\norm{\hphantom O}_1)$ kontinu. Kedua topologi sama tepat ketika
identitas sebaliknya kontinu, yaitu ketika ada $c>0$ sehingga
\[
 \norm x\leq c\norm{x}_1\quad(x\in H). \tag{1}
\]
Jika (1) berlaku dan $\delta=c^{-2}$, maka
\[
 \delta\norm x^2\leq\langle Tx,x\rangle
 \leq\norm{Tx}\norm x,
\]
sehingga $\norm{Tx}\geq\delta\norm x$. Jadi $T$ injektif dan jangkauannya
tertutup. Selain itu
$(\ran T)^\perp=\ker T^*=\ker T=\{0\}$, sehingga jangkauannya padat. Ia
karena itu sama dengan~$H$, dan invers $T^{-1}$ terbatas dengan norma paling
besar $\delta^{-1}$.

Sebaliknya, andaikan $T$ invertibel. Inversnya positif: jika $y=Tx$, maka
\[
 \langle T^{-1}y,y\rangle=\langle x,Tx\rangle
 =\langle Tx,x\rangle\geq0.
\]
Gunakan Cauchy--Schwarz untuk bentuk positif $[u,v]=\langle Tu,v\rangle$:
\begin{align*}
 \norm x^2
 &=\langle T(T^{-1}x),x\rangle\\
 &\leq
 \langle T(T^{-1}x),T^{-1}x\rangle^{1/2}
 \langle Tx,x\rangle^{1/2}\\
 &=\langle x,T^{-1}x\rangle^{1/2}\norm{x}_1
 \leq\norm{T^{-1}}^{1/2}\norm x\,\norm{x}_1.
\end{align*}
Untuk $x\neq0$, pembagian dengan $\norm x$ memberi
$\norm x\leq\norm{T^{-1}}^{1/2}\norm{x}_1$, yaitu (1). Ini menyelesaikan
(b).

Untuk (c), tetapkan $A=T^{-1}S^*T$, yang terbatas karena ketiga faktornya
terbatas. Untuk semua $x,y\in H$,
\begin{align*}
 \langle Sx,y\rangle_1
 &=\langle TSx,y\rangle
  =\langle Sx,Ty\rangle
  =\langle x,S^*Ty\rangle\\
 &=\langle Tx,T^{-1}S^*Ty\rangle
  =\langle x,Ay\rangle_1.
\end{align*}
Jadi, menurut definisi adjoin dalam hasil kali dalam baru,
$S^{*_1}=A=T^{-1}S^*T$. Urutan faktor ini juga memeriksa bahwa rumus sesuai
dengan konvensi edisi bahwa hasil kali dalam linear pada variabel pertama.
\end{o001proof}
\end{o001solution}
